\section{Method}
\label{sec:method}

% ============================================================
% OUTLINE -- Method            target: ~1.5-1.7 pages (the core)
% Spine: setup -> scaffolding (bounded cost) -> dense per-round
% target (telescoping => supervised, not RL) -> the loss.
%
% Notation (use consistently paper-wide):
%   G = (V,E)  code graph; edges = contains/call/import/co-change
%   Y*         gold set of code units that must change
%   L          the fixed-length returned list; recall@L the metric
%   K          window size: ranked candidates revealed per round
%   P_r        per-round candidate pool
%   Phi_r      secured-gold potential
% Display equations in the body: the telescoping identity, the
% weighted BCE, the coverage term. Three equations max.
% ============================================================

\subsection{Problem Setup}
\label{sec:setup}
% - Repository R, natural-language issue q. A code graph
%   G = (V,E): nodes are code units (file / class / function),
%   edges are structural relations (contains, call, import,
%   co-change).
% - Task: recover the gold set Y* of units that must be modified.
% - Metric: recall@L over a fixed-length returned list L.
% - Frame the task as SET RECOVERY, not ranking -- this framing
%   is exactly what the objective in 3.4 is built to match.
%   (~0.2-0.25 page)

\subsection{Multi-Round Localization Scaffolding}
\label{sec:scaffold}

The scaffolding processes a single issue $q$ over $R$ rounds. A
retriever (dense, lexical, or their union) returns a ranked
list of code units, and each round consumes the next $K$ unseen
entries of that list. At the start of round $r$, the
scaffolding holds three pieces of state: the current returned
list $L_{r-1}$ (initialised to $\emptyset$ at $r{=}0$), the set
of node identifiers $\mathcal{S}_{r-1}$ already shown to the
model, and an offset $rK$ into the ranked list. The loop
terminates when the ranked list is exhausted or a maximum
number of rounds $R_{\max}$ is reached. An optional
confident-stop hook is left for the RL stage
(Section~\ref{sec:conclusion}); the supervised training of
Section~\ref{sec:objective} does not predict it, so a
trajectory always runs to $R_{\max}$ or exhaustion.

Each round combines a bounded window of newly-ranked candidates
with a bounded graph-neighborhood expansion, then refines $L$;
Figure~\ref{fig:scaffold} sketches the loop.

\begin{figure}[t]
\centering
% TODO: replace with the multi-round scaffolding diagram.
\fbox{\rule{0pt}{2.4cm}\rule{0.9\columnwidth}{0pt}}
\caption{The multi-round localization scaffolding. Each round
consumes a bounded window of $K$ ranked candidates, expands
neighborhoods around $C$ seed centers via bounded BFS on the
code graph, and refines a fixed-length returned list $L_r$.
Per-call token cost stays roughly constant across rounds.}
\label{fig:scaffold}
\end{figure}

\paragraph{Window.}
The scaffolding pulls the next $K$ unseen entries of the ranked
list starting at offset $rK$, skipping any node already in
$\mathcal{S}_{r-1}$. We denote this window $W_r$. If $W_r$ is
empty the round terminates the loop.

\paragraph{Step A: seed-center selection.}
A single model call selects $C$ seed centers from $W_r$, i.e.\
nodes whose graph neighborhoods are most likely to contain
unseen gold units. If the model returns fewer than $C$ valid
identifiers we pad with the top-ranked nodes of $W_r$ not yet
returned, so Step~A always emits exactly $C$ centers.

\paragraph{Step B: neighborhood expansion.}
For each of the $C$ centers, the scaffolding runs a depth-$d$
BFS along structural edges of the code graph (\textit{contains}
by default; ablations in Section~\ref{sec:ablation} cover the
four edge types). The neighbors are filtered to the next $N$
unseen entries of the ranked list, so the candidate set the
model sees stays bounded. A model call per center then selects
the subset of these neighbors that must be modified, and the
union across centers, $\mathcal{N}_r$, is deduplicated.

\paragraph{Step C: list refinement.}
A final model call sees the issue $q$, the current list
$L_{r-1}$, the window $W_r$, and the deduplicated set
$\mathcal{N}_r$, and selects $L$ identifiers from
$L_{r-1} \cup W_r \cup \mathcal{N}_r$ to form $L_r$. The
scaffolding then records the newly-shown identifiers, advances
the offset, and continues.

Three properties of this loop matter for what follows.
\emph{First}, every model call sees only $W_r$, a single
neighborhood, or the deduplicated round pool, so the per-call
token cost is roughly constant across rounds; a 32K context
comfortably absorbs each call regardless of how many rounds the
trajectory uses. \emph{Second}, the loop offers a natural
ablation: removing Steps~A and~B reduces it to a pure
window-refine loop and isolates the contribution of graph
exploration. \emph{Third}, $L_r$ at any round is a valid
prediction, so the round structure exposes a per-round
potential that we develop into a training target in the next
section.

\subsection{A Dense, Recall-Aligned Per-Round Target}
\label{sec:telescope}
% - Define the secured-gold potential  Phi_r = |Y* intersect L_r|.
% - Per-round signal  r_r = (Phi_r - Phi_{r-1}) / D, with
%   D = |Y* intersect candidates|  =  (gold added - gold dropped)/D.
% - EQUATION (telescoping):  sum_r r_r = (Phi_R - Phi_0)/D
%   = recall@L(L_final) - recall@L(L_0).
% - THEOREM (state here, proof -> Appendix A): by potential-
%   based reward shaping \citep{ng1999policy}, shaping with a
%   state potential leaves the optimal policy invariant; hence
%   the dense per-round objective and sparse end-to-end
%   recall@L share the same optimum -- there is no myopia gap.
% - IMPLICATION (the load-bearing sentence): the per-round
%   optimum is closed-form, so training needs no policy
%   gradient; supervised learning on the per-round target is
%   denser, more stable, and cheaper than RL on the same signal.
%   (~0.3-0.35 page)

\subsection{Supervised Training Objective}
\label{sec:objective}
% - Per round build the pool P_r = L_{r-1} U window_r; labels
%   y_j = 1[j in Y*].
% - Scoring without a new head: use the model's own yes/no
%   token logits, s_j = logit(yes) - logit(no), p_j = sigma(s_j).
%   Zero new parameters; LoRA-friendly; avoids the cold-start
%   failure of a randomly-initialized head.
% - EQUATION (weighted BCE): the main term; its optimum is
%   "every gold scored above every non-gold", which yields a
%   max-coverage top-L.
% - EQUATION (coverage term L_cov): a differentiable, listwise,
%   permutation-invariant surrogate for covering the whole gold
%   set; it never penalizes gold-vs-gold order.
% - Total objective: L = L_bce + lambda * L_cov.
% - State the guarantee scope: exact recall@L co-optimality
%   holds at w+ = w- = 1; an asymmetric w+ > w- is a class-
%   imbalance / "dropping gold is worse" heuristic -- report it
%   as a tuning choice, not part of the theorem. lambda = 0 is
%   the ablation baseline.
% - Top-L selection is inference-only and never differentiated.
% - One sentence: stages are trained densest-signal-first
%   (curriculum); full staging details -> Appendix B.
%   (~0.45-0.5 page)
